% !TeX program = lualatex-dev
% ===================================================================
% mwe_leg_headings.tex — MWE-1 da camada Leg (crpsp-leg.sty)
%
% Valida os AGRUPADORES como headings tagueados fora do TOC:
%   - \LegAgrup em 4 profundidades sob \chapter (H1) + \section (H2)
%     deve produzir <H3>…<H6> na árvore de tags;
%   - o TOC não pode listar nenhum agrupador;
%   - contadores e marks do book intactos (seção seguinte numerada
%     corretamente).
%
% Compilar 2x: lualatex-dev mwe_leg_headings.tex
% Inspecionar: árvore de tags (Acrobat/PAC) ou \showtags no log.
% ===================================================================

\DocumentMetadata{
    lang        = pt-BR,
    pdfstandard = ua-2,
    pdfversion  = 2.0,
    testphase   = {phase-III,table,firstaid},
    tagging     = on,
}

\documentclass[11pt,a5paper,twoside,openany]{book}

\usepackage{../../book-crpsp_acessivel}

% dc:title no XMP — exigência PDF/UA-2 (veraPDF 8.11.1-t1)
\hypersetup{pdftitle={MWE Leg: agrupadores como headings tagueados}}

\begin{document}

\TOCAcessivel
\mainmatter

\chapter{Capítulo de teste}          % H1

Parágrafo introdutório do capítulo.

\section{Norma transcrita}           % H2

Parágrafo antes da norma. Os agrupadores abaixo devem virar
H3–H6 sem aparecer no sumário.

\begin{LegNorma}{urn:lex:br:teste:lei:2026;1}{LEI DE TESTE Nº 1, DE 2026}

\LegEmenta{Ementa de teste para o MWE de headings.}

\LegAgrup[1]{livro}{LIVRO I}{PARTE GERAL}          % esperado: H3

Texto sob o livro.

\LegAgrup[2]{titulo}{TÍTULO I}{DISPOSIÇÕES PRELIMINARES} % esperado: H4

Texto sob o título.

\LegAgrup[3]{capitulo}{CAPÍTULO I}{DISPOSIÇÕES GERAIS}   % esperado: H5

Texto sob o capítulo.

\LegAgrup[4]{secao}{SEÇÃO I}{DO OBJETO}                  % esperado: H6

\LegArtigo{Art. 1º}{Artigo plano sob os agrupadores, para conferir
que a lista continua tagueada normalmente após os headings.}

\end{LegNorma}

Parágrafo após a norma: fonte e espaçamento devem ter voltado
ao normal (serifada, 1,5 linha).

\section{Seção seguinte}             % deve numerar 1.2 (contador intacto)

O número desta seção confere que os agrupadores não incrementaram
o contador de seções.

\end{document}
